%%=============================================================================
%% Methodologie
%%=============================================================================

\chapter{\IfLanguageName{dutch}{Methodologie}{Methodology}}%
\label{ch:methodologie}

Dit hoofdstuk beschrijft de onderzoeksaanpak en het plan van aanpak van deze bachelorproef. Waar Hoofdstuk~\ref{ch:stand-van-zaken} het theoretisch kader rond ERP-implementaties, KMO-context en e-invoicing uiteenzette, wordt hier toegelicht hoe deze inzichten werden vertaald naar een praktijkgericht onderzoek binnen de casus van Hanuman. 

Dit hoofdstuk blijft bewust op het niveau van onderzoeksopzet en fasering. De inhoudelijke uitwerking van de implementatie, configuratiekeuzes en evaluatieresultaten wordt in de daaropvolgende hoofdstukken behandeld.

\section{Onderzoeksstrategie}

\subsection{Casestudybenadering}

Er werd gekozen voor een praktijkgerichte enkelvoudige casestudy. Deze methode is geschikt wanneer een fenomeen wordt onderzocht binnen zijn reële context en wanneer organisatorische en technische factoren sterk met elkaar verweven zijn. ERP-implementaties binnen KMO’s voldoen aan deze kenmerken.

Hanuman fungeert als representatieve casus van een servicegerichte Belgische KMO met beperkte personeelsomvang, hybride processen (verkoop en herstellingen) en een externe verplichting tot gestructureerde e-facturatie via Peppol.

Het doel is analytische generaliseerbaarheid: het ontwikkelen van een onderbouwd migratiekader dat toepasbaar is op gelijkaardige organisaties.

\subsection{Ontwerpgericht karakter van het onderzoek}

Het onderzoek heeft een ontwerpgericht karakter, aangezien het niet enkel een probleem analyseert, maar ook een oplossing ontwerpt en evalueert. Dit resulteert in:

\begin{itemize}
    \item een gefaseerd migratiekader,
    \item een geconfigureerde proof-of-concept in Odoo,
    \item een evaluatiekader op basis van proces-KPI’s.
\end{itemize}

Hierdoor wordt een brug geslagen tussen theoretische inzichten en praktische toepassing.

\section{Gefaseerde onderzoeksopzet}

Het onderzoek werd uitgevoerd in zes opeenvolgende fasen. Voor elke fase worden doelstelling, methode en deliverables kort toegelicht.

\subsection{Fase 0: Literatuurstudie}

\textbf{Doelstelling.}  
Opbouwen van een theoretisch kader rond ERP, implementatierisico’s, KMO-specifieke uitdagingen, procesherontwerp en e-invoicing.

\textbf{Methode.}  
Analyse en synthese van wetenschappelijke literatuur en officiële beleidsdocumenten.

\textbf{Deliverable.}  
Het theoretisch kader zoals uitgewerkt in Hoofdstuk~\ref{ch:stand-van-zaken}.

\subsection{Fase 1: Analyse van de huidige situatie (AS-IS)}

\textbf{Doelstelling.}  
In kaart brengen van bestaande processen en identificeren van inefficiënties, foutbronnen en redundantie.

\textbf{Onderzoeksmethoden.}
\begin{itemize}
    \item Observatie van dagelijkse administratieve workflows
    \item Informele semigestructureerde interviews met stakeholders
    \item Analyse van bestaande Excel-bestanden en documenten
\end{itemize}

\textbf{Deliverables.}
\begin{itemize}
    \item AS-IS procesbeschrijving
    \item Overzicht van knelpunten en risico’s
\end{itemize}

\subsection{Fase 2: Requirements-analyse en scopebepaling}

\textbf{Doelstelling.}  
Vertalen van vastgestelde knelpunten naar functionele en niet-functionele vereisten voor de ERP-oplossing.

\textbf{Onderzoeksmethoden.}
\begin{itemize}
    \item Requirements-elicitation via stakeholderafstemming
    \item Mapping van bedrijfsprocessen naar ERP-standaardfunctionaliteiten
\end{itemize}

\textbf{Deliverables.}
\begin{itemize}
    \item Functioneel vereistenoverzicht
    \item Afgebakende implementatiescope
\end{itemize}

\subsection{Fase 3: Ontwerp van de gewenste situatie (TO-BE) en migratiekader}

\textbf{Doelstelling.}  
Ontwerpen van een toekomstige processtructuur en bepalen van een gefaseerde implementatiestrategie.

\textbf{Onderzoeksmethoden.}
\begin{itemize}
    \item Procesherontwerp op basis van ERP-standaardlogica
    \item Risicoanalyse geïnspireerd door literatuur rond implementatiesuccesfactoren
\end{itemize}

\textbf{Deliverables.}
\begin{itemize}
    \item TO-BE procesmodel
    \item Gefaseerd migratiekader
\end{itemize}

\subsection{Fase 4: Proof-of-Concept configuratie}

\textbf{Doelstelling.}  
Praktische validatie van de ontworpen processen in een testomgeving.

\textbf{Onderzoeksmethoden.}
\begin{itemize}
    \item Configuratie van relevante Odoo-modules
    \item Invoer van representatieve stamdata
    \item Scenario-gebaseerde functionele tests
\end{itemize}

\textbf{Deliverables.}
\begin{itemize}
    \item Geconfigureerde testomgeving
    \item Gevalideerde procesflows
\end{itemize}

\subsection{Fase 5: Evaluatie via KPI’s en e-invoicing-voorbereiding}

\textbf{Doelstelling.}  
Objectiveren van procesverbetering en beoordelen van voorbereiding op gestructureerde e-facturatie.

\textbf{Onderzoeksmethoden.}
\begin{itemize}
    \item Operationalisatie van proces-KPI’s
    \item Vergelijking tussen AS-IS analyse en gesimuleerde TO-BE scenario’s
    \item Analyse van configuratievereisten voor e-invoicing
\end{itemize}

\textbf{Deliverables.}
\begin{itemize}
    \item KPI-evaluatiekader
    \item Overzicht van aandachtspunten voor Peppol-compliance
\end{itemize}

\section{Kwaliteitsbewaking en beperkingen}

De betrouwbaarheid van het onderzoek wordt verhoogd door triangulatie: observaties worden gecombineerd met stakeholderinput en documentanalyse. De evaluatie is gedeeltelijk gebaseerd op gesimuleerde scenario’s, aangezien historische logdata in de oorspronkelijke situatie beperkt beschikbaar is.

De resultaten zijn voornamelijk toepasbaar op gelijkaardige servicegerichte KMO’s en beogen analytische in plaats van statistische generaliseerbaarheid.

\section{Samenvatting}

De methodologie bestaat uit een gefaseerde aanpak met analyse, ontwerp, proof-of-concept en evaluatie. Dit hoofdstuk beschreef het plan van aanpak en de gekozen onderzoeksmethoden. In de volgende hoofdstukken worden de uitvoering en resultaten van deze fasen inhoudelijk besproken, waarna in Hoofdstuk~\ref{ch:conclusie} een antwoord wordt geformuleerd op de onderzoeksvragen.